%% ============================================================
%% quasar/horizon-soundness.tex
%% ============================================================
%%
%% Canonical Horizon (Prism-bound) soundness theorem.
%% ============================================================

\subsection*{Horizon soundness}

\begin{theorem}[Horizon-final implies all lanes verify the same transcript]
\label{thm:horizon-soundness}
Let $\mathsf{cert}$ be a Horizon certificate
(Definition~\ref{def:horizon-cert}) with
$\mathsf{Prism}(\mathsf{cert}) = 1$. Then the underlying validator set
of size $n$ with $f < n/3$ corruptions cannot have produced a Horizon-
final certificate for two distinct $\mathsf{NebulaRoot}$ values at the
same epoch and round, except with negligible probability under:
\begin{enumerate}[nosep]
  \item BLS12-381 co-CDH (for the Beam lane);
  \item Module-LWE / Module-SIS (for the Pulsar lane,
        Theorem~\ref{thm:pulsar-suf-cma});
  \item EUF-CMA security of ML-DSA-65 / FIPS 204 (for the ML-DSA cert
        set lane).
\end{enumerate}
The proof is by independent reduction: any pair of Prism-accepted certs
on conflicting bundle roots forces a forgery in at least one lane.
This is independent across the three lanes (different hardness
assumptions); a quantum adversary breaking BLS does not break Pulsar
or ML-DSA, and vice-versa.

\textbf{Companion proof:} \texttt{quasar-cert-soundness.tex} Theorem
7.5 (classical soundness), 7.6 (parallel-lane liveness), 7.7 (PQ
safety).
\end{theorem}

\begin{remark}[Groth16 wrapper does not contribute PQ security]
\label{rem:groth16-wrapper}
A Z-Chain Groth16 rollup proof of ML-DSA cert-set verification
(Definition~\ref{def:proof-lane}) is a \emph{classical succinct proof of
post-quantum signature verification}. The proof system is
pairing-based and is not post-quantum. Horizon's PQ soundness
(Theorem~\ref{thm:horizon-soundness} parts 2 and 3) flows from the
ML-DSA signatures and Pulsar threshold signature alone; the Groth16
proof is verified for compatibility / compression only and is not part
of the PQ reduction. See Remark~\ref{rem:groth16-not-pq}.
\end{remark}
